\documentclass[11pt,a4paper]{article}
\usepackage[utf8]{inputenc}
\usepackage[T1]{fontenc}
\usepackage[spanish,es-noshorthands]{babel}
\usepackage{geometry}
\geometry{margin=2.6cm}
\usepackage{calc}
\usepackage{longtable}
\usepackage{booktabs}
\usepackage{array}
\usepackage{ragged2e}
\usepackage{graphicx}
\usepackage{float}
\usepackage[table]{xcolor}
\usepackage{caption}
\captionsetup{skip=6pt}
\usepackage{titlesec}
\usepackage{fancyhdr}
\usepackage[hidelinks]{hyperref}
\usepackage{enumitem}
\usepackage{parskip}
\usepackage{needspace}
\usepackage{etoolbox}
\setlength{\parskip}{0.55em}
\setlist{itemsep=0.2em}

\pretocmd{\section}{\needspace{5\baselineskip}}{}{}
\pretocmd{\subsection}{\needspace{4\baselineskip}}{}{}
\pretocmd{\subsubsection}{\needspace{3\baselineskip}}{}{}

\providecommand{\tightlist}{%
  \setlength{\itemsep}{0pt}\setlength{\parskip}{0pt}}
\providecommand{\pandocbounded}[1]{#1}
\newcommand{\passthrough}[1]{#1}

\hypersetup{
  pdftitle={HemoScan Andino --- Arquitectura Integrada},
  pdfauthor={Equipo HemoScan Andino -- UPeU Juliaca},
  colorlinks=false,
  bookmarksopen=true
}

\setcounter{secnumdepth}{-2}
\setcounter{tocdepth}{3}
\titleformat{\section}{\normalfont\Large\bfseries}{}{0em}{}
\titleformat{\subsection}{\normalfont\large\bfseries}{}{0em}{}
\titleformat{\subsubsection}{\normalfont\normalsize\bfseries}{}{0em}{}

\pagestyle{fancy}
\fancyhf{}
\fancyhead[L]{\small HemoScan Andino}
\fancyhead[R]{\small Arquitectura Integrada --- CE01 \textperiodcentered{} CE02 \textperiodcentered{} CE03}
\fancyfoot[C]{\thepage}
\renewcommand{\headrulewidth}{0.4pt}

% figura de diagrama (imagen vectorial renderizada con mermaid-cli)
\newcommand{\diag}[3]{%
  \begin{figure}[H]\centering
  \includegraphics[width=#2\linewidth,height=0.86\textheight,keepaspectratio]{assets/11-arquitectura/#1}%
  \caption*{#3}\end{figure}}

\begin{document}

\begin{titlepage}
\centering
\vspace*{1.2cm}
{\large UNIVERSIDAD PERUANA UNI\'ON\par}
{\normalsize Facultad de Ingenier\'ia y Arquitectura\par}
{\normalsize Escuela Profesional de Ingenier\'ia de Sistemas --- Filial Juliaca\par}
\vspace{0.5cm}
\rule{0.6\linewidth}{0.4pt}
\vspace{1.6cm}

{\large PROYECTO DE TESIS / PROYECTO INTEGRADOR\par}
\vspace{1.0cm}
{\huge\bfseries HemoScan Andino\par}
\vspace{0.6cm}
{\Large Sistema no invasivo de tamizaje de anemia infantil con fusi\'on \'optica multisensor y autocalibraci\'on por altitud\par}
\vspace{1.6cm}
\rule{0.6\linewidth}{0.4pt}
\vspace{1.6cm}

{\Large\bfseries Arquitectura Integrada del Sistema\par}
\vspace{0.5cm}
{\normalsize Documento transversal que unifica las tres \'areas de competencia:\par}
{\normalsize CE01 (Gesti\'on de TI) + CE02 (Ingenier\'ia de Software) + CE03 (Infraestructura Tecnol\'ogica)\par}
\vspace{2.0cm}

{\normalsize Juliaca, Puno, Per\'u \par}
{\normalsize Julio de 2026\par}
\vfill
\end{titlepage}

\tableofcontents
\newpage

\textbf{Documento:} Arquitectura Integrada de HemoScan Andino (vista transversal del proyecto)

\textbf{Prop\'osito:} consolidar en un \'unico modelo arquitect\'onico coherente y trazable las decisiones de las tres \'areas de competencia del proyecto, con \'enfasis en la arquitectura del sistema (vistas, componentes, decisiones y su justificaci\'on).

\textbf{Sistema:} HemoScan Andino --- dispositivo \'optico multisensor (ESP32-S3), aplicaci\'on m\'ovil Flutter/PWA \emph{offline-first} con inferencia embebida y autocalibraci\'on por altitud, backend de interoperabilidad FHIR (FastAPI, monolito modular), base de datos PostgreSQL 16 + PostGIS 3.4 y plataforma web de anal\'itica geoespacial (React/TypeScript).

\textbf{Organizaci\'on (caso ilustrativo):} Microred de Salud Altiplano-Puno.

{\small
\begin{longtable}[]{@{}
  >{\raggedright\arraybackslash}p{(\linewidth - 4\tabcolsep) * \real{0.08}}
  >{\raggedright\arraybackslash}p{(\linewidth - 4\tabcolsep) * \real{0.52}}
  >{\raggedright\arraybackslash}p{(\linewidth - 4\tabcolsep) * \real{0.40}}@{}}
\toprule\noalign{}
N.° & Integrante & Rol / \'Area de competencia \\
\midrule\noalign{}
\endhead
\bottomrule\noalign{}
\endlastfoot
1 & \textbf{{[}DATO PENDIENTE DE VALIDAR: Integrante 1{]}} & CE03 --- Infraestructura Tecnol\'ogica (red, seguridad, centro de datos) \\
2 & \textbf{{[}DATO PENDIENTE DE VALIDAR: Integrante 2{]}} & CE01 --- Gesti\'on de TI (diagn\'ostico, business case, plan, procesos) \\
3 & \textbf{{[}DATO PENDIENTE DE VALIDAR: Integrante 3{]}} & CE02 --- Ingenier\'ia de Software (requerimientos, datos, sistema, calidad) \\
\end{longtable}
}

\textbf{Lugar y fecha:} Juliaca, Puno, Per\'u --- 07 de julio de 2026.

\begin{quote}
\textbf{Nota metodol\'ogica.} Este documento es una vista transversal que \emph{integra} y no reemplaza a los nueve entregables oficiales del proyecto. Todos los diagramas se especifican como \emph{diagrama-como-c\'odigo} (notaci\'on Mermaid) y se renderizan a imagen vectorial con un renderizador libre de licencia (\texttt{mermaid-cli}); las im\'agenes incluidas son la salida directa de ese renderizado. Todo valor o supuesto que no se derive de los entregables previos se marca expl\'icitamente como \textbf{{[}DATO PENDIENTE DE VALIDAR{]}}. La arquitectura respeta, sin contradecirlas, las decisiones (ADR) y restricciones ya establecidas en los Entregables N.° 1 a 9.
\end{quote}

\begin{center}\rule{0.5\linewidth}{0.5pt}\end{center}

\section{1. Resumen ejecutivo de la arquitectura}\label{resumen}

HemoScan Andino se concibe como un sistema socio-t\'ecnico de tamizaje ---no de diagn\'ostico aut\'onomo--- desplegado en el primer nivel de atenci\'on de salud en un contexto altoandino con conectividad y energ\'ia intermitentes. Su arquitectura est\'a gobernada por cuatro \emph{drivers} que condicionan casi todas las decisiones: \textbf{(i) operaci\'on offline-first} (la funci\'on cl\'inica cr\'itica no puede depender de conectividad permanente); \textbf{(ii) auditabilidad y supervisi\'on humana} exigidas por el DS 115-2025-PCM y la Ley N.° 29733; \textbf{(iii) interoperabilidad est\'andar} con el ecosistema del MINSA mediante HL7 FHIR; y \textbf{(iv) bajo costo}, coherente con el presupuesto de Fase 0 (~S/ 320 por 8 meses).

De estos drivers se derivan las decisiones estructurales centrales, ya formalizadas como ADR en entregables previos y reafirmadas aqu\'i: una arquitectura \textbf{offline-first con patr\'on \emph{outbox}} (ADR-01); un \textbf{monolito modular} de siete paquetes de dominio en el backend (ADR-02, ADR-03), evitando la complejidad operativa de microservicios en Fase 0 sin renunciar a l\'imites de dominio claros; el \textbf{desacoplamiento} entre el motor de inferencia y el m\'odulo de autocalibraci\'on por altitud (ADR-05), de modo que la f\'ormula de correcci\'on pueda evolucionar sin reentrenar el modelo; una \textbf{base de datos \'unica} PostgreSQL + PostGIS (ADR-10) con cuatro esquemas de dominio; \textbf{Keycloak} como proveedor de identidad OIDC/RBAC (ADR-09); y el \textbf{aprendizaje federado empaquetado pero inactivo} en Fase 0 (ADR-08), evitando deuda t\'ecnica futura sin sobredimensionar el piloto.

La principal contribuci\'on de esta vista es demostrar que las tres \'areas de competencia del proyecto convergen en un \'unico modelo: la Gesti\'on de TI (CE01) aporta los \emph{drivers}, procesos y el alcance por fases; la Ingenier\'ia de Software (CE02) aporta los componentes, el modelo de datos y las interfaces; y la Infraestructura Tecnol\'ogica (CE03) aporta el despliegue, la red, la seguridad y el dimensionamiento de capacidad. La secci\'on 12 formaliza esa convergencia con una matriz de trazabilidad.

\section{2. Atributos de calidad y escenarios}\label{qas}

La arquitectura se justifica frente a atributos de calidad expresados como escenarios verificables (est\'imulo → respuesta → m\'etrica), siguiendo la disciplina de \emph{Quality Attribute Scenarios}. Las m\'etricas marcadas como supuestas requieren validaci\'on en banco de pruebas.

{\small
\begin{longtable}[]{@{}
  >{\raggedright\arraybackslash}p{(\linewidth - 6\tabcolsep) * \real{0.16}}
  >{\raggedright\arraybackslash}p{(\linewidth - 6\tabcolsep) * \real{0.50}}
  >{\raggedright\arraybackslash}p{(\linewidth - 6\tabcolsep) * \real{0.34}}@{}}
\toprule\noalign{}
Atributo & Escenario (est\'imulo → respuesta) & M\'etrica objetivo \\
\midrule\noalign{}
\endhead
\bottomrule\noalign{}
\endlastfoot
Disponibilidad (offline) & Un puesto de salud unipersonal pierde conectividad durante una jornada; el personal sigue tamizando y registrando. & Funci\'on cl\'inica cr\'itica operativa \textbf{72 h} sin conectividad; $\geq$ 50 capturas en buffer local. \\
Rendimiento & Tras capturar las tres se\~nales \'opticas, la app ejecuta la inferencia embebida. & Resultado de tamizaje en \textbf{$<$ 3 s} en gama baja-media. \\
Interoperabilidad & Al recuperar conectividad, la app sincroniza recursos FHIR con el HIS-MINSA. & 100\% de recursos como \textbf{FHIR v\'alido}; sincronizaci\'on idempotente por UUID. \\
Auditabilidad & Un auditor consulta el origen de una determinaci\'on cl\'inica. & Traza completa \textbf{Provenance + AuditEvent} por cada Observation, sin excepci\'on (RN-11). \\
Privacidad & Un actor sin autorizaci\'on intenta leer datos a nivel individual desde la web. & Bloqueo por RBAC + RLS; la capa geoespacial \textbf{no} tiene FK a paciente (RN-08). \\
Seguridad & Se compromete una credencial de aplicaci\'on de rol auditor. & Acceso limitado a solo lectura por RLS; cifrado por columna protege PII; evento en \texttt{pgAudit}. \\
Mantenibilidad & Cambia la f\'ormula de correcci\'on por altitud (NTS 213). & Se versiona en \texttt{config} sin reentrenar el modelo ni tocar el motor de inferencia (ADR-05). \\
Portabilidad / costo & El piloto debe correr en un \'unico servidor de bajo costo. & Backend + BD + Keycloak en \textbf{1 instancia}; presupuesto Fase 0 ~S/ 320 / 8 meses. \\
\end{longtable}
}

\section{3. Vista de contexto (C4 --- nivel 1)}\label{contexto}

La vista de contexto delimita el sistema frente a sus actores humanos (cuidadores, personal CRED, personal m\'edico, jefatura de microred, analista DIRESA/GERESA, administrador y auditor) y a los sistemas externos del ecosistema de salud p\'ublico (HIS-MINSA, Padr\'on Nominal, SIGA y NETLAB-INS). El principio rector ---tamizaje y nunca diagn\'ostico aut\'onomo--- se materializa ya en este nivel: el personal m\'edico es un actor de decisi\'on obligatorio, no un consumidor pasivo de resultados.

\diag{d01-contexto.pdf}{0.95}{Diagrama 1 --- Contexto del sistema (C4 nivel 1): actores y sistemas externos.}

\section{4. Vista de contenedores (C4 --- nivel 2)}\label{contenedores}

Los \emph{contenedores} (unidades desplegables independientes) son cinco: el nodo/dispositivo con su firmware, la aplicaci\'on m\'ovil, el backend de interoperabilidad, la plataforma web y la base de datos, m\'as Keycloak como servicio de identidad. La comunicaci\'on dispositivo-app es local (BLE GATT); toda comunicaci\'on con el backend es REST sobre HTTPS con carga \'util FHIR y, desde la app, \textbf{diferida} (mediada por la cola de \emph{outbox}). Ning\'un contenedor mantiene su propio directorio de usuarios: la identidad se delega \'integramente a Keycloak (ADR-09).

\diag{d02-contenedores.pdf}{0.95}{Diagrama 2 --- Contenedores del sistema (C4 nivel 2) y protocolos entre ellos.}

\section{5. Vista de componentes (C4 --- nivel 3)}\label{componentes}

\subsection{5.1 Backend de interoperabilidad}\label{comp-backend}

El backend es un \textbf{monolito modular} (ADR-02): un \'unico proceso desplegable organizado en siete paquetes de dominio con l\'imites expl\'icitos. Todos dependen del paquete \texttt{identidad} para autorizaci\'on; el paquete \texttt{clinico\_fhir} es el n\'ucleo del que penden notificaciones, consejer\'ia y anal\'itica geoespacial. El paquete \texttt{aprendizaje\_federado} existe como componente desde Fase 0 pero permanece inactivo (ADR-08), evitando la deuda de una reintroducci\'on futura.

\diag{d03-componentes-backend.pdf}{0.9}{Diagrama 3 --- Componentes del backend modular (siete paquetes de dominio).}

\subsection{5.2 Aplicaci\'on m\'ovil offline-first}\label{comp-app}

La app concentra la l\'ogica cl\'inica cr\'itica en el dispositivo del personal de salud. El adaptador BLE captura de forma sincronizada las tres se\~nales \'opticas; el motor de inferencia embebido (ONNX Runtime Mobile, ruta alterna TFLite) estima la hemoglobina; el m\'odulo de autocalibraci\'on la corrige consultando un Modelo Digital de Elevaci\'on \textbf{offline} a partir del GPS (nunca altitud cruda del GPS); y la cola de \emph{outbox} garantiza persistencia local at\'omica y sincronizaci\'on diferida e idempotente.

\diag{d04-componentes-app.pdf}{0.85}{Diagrama 4 --- Componentes de la aplicaci\'on m\'ovil offline-first.}

\section{6. Vista de ejecuci\'on (runtime)}\label{runtime}

\subsection{6.1 Flujo cr\'itico de tamizaje}\label{flujo-critico}

El escenario de ejecuci\'on m\'as representativo recorre el sistema de extremo a extremo: captura multisensor, inferencia embebida, correcci\'on por altitud auditable, clasificaci\'on de severidad, persistencia at\'omica de los recursos FHIR y sincronizaci\'on diferida. En ninguna trayectoria el sistema confirma o descarta anemia de forma aut\'onoma: la severidad bajo umbral genera una alerta de derivaci\'on que exige confirmaci\'on humana (RN-01, DS 115-2025-PCM).

\diag{d06-secuencia-critica.pdf}{1.0}{Diagrama 5 --- Secuencia del flujo cr\'itico de tamizaje (captura → inferencia → autocalibraci\'on → outbox → sincronizaci\'on).}

\subsection{6.2 Ciclo de vida de la sincronizaci\'on offline}\label{outbox}

La contraparte de estados del flujo anterior es el ciclo de vida de cada registro en la cola de \emph{outbox} (ADR-01): generado localmente, encolado, transmitido, y sincronizado o marcado como conflicto pendiente de revisi\'on manual. El backoff exponencial y la idempotencia por UUID hacen segura la reintenci\'on ante conectividad intermitente.

\diag{d07-estados-outbox.pdf}{0.9}{Diagrama 6 --- M\'aquina de estados de la sincronizaci\'on offline (patr\'on outbox).}

\section{7. Vista de datos}\label{datos}

La arquitectura de datos reside en una base \'unica PostgreSQL 16 + PostGIS 3.4 (ADR-10) con cuatro esquemas de dominio (\texttt{clinico}, \texttt{geo}, \texttt{config}, \texttt{gestion}) que reproducen a nivel de base de datos la misma modularidad del backend. Los seis recursos HL7 FHIR del proyecto (Patient, Consent, Device, Observation, Provenance, AuditEvent) se serializan a partir de esas tablas; la referencia polim\'orfica de \texttt{AuditEvent} se resuelve con una tabla t\'ecnica compartida (\texttt{recurso\_fhir}) que preserva integridad referencial real, y la capa geoespacial agregada (hex\'agonos H3) se dise\~na \textbf{sin} clave for\'anea a paciente para materializar la privacidad por dise\~no (RN-08).

\diag{d05-datos-fhir.pdf}{0.95}{Diagrama 7 --- Arquitectura de datos: esquemas de dominio y mapeo a los seis recursos FHIR.}

\section{8. Vista de despliegue e infraestructura (CE03)}\label{despliegue}

\subsection{8.1 Despliegue por fases}\label{deploy-fases}

En Fase 0, un \textbf{\'unico servidor} de bajo costo aloja el backend, la base de datos y Keycloak, consistente con el presupuesto comprometido; el nodo dedicado de aprendizaje federado se documenta como extensi\'on de Fase 1/2 pero no se despliega. La validaci\'on de la capacidad real de ese servidor \'unico es responsabilidad del paquete EDT 2.2 de Infraestructura (CE03).

\diag{d08-despliegue-fases.pdf}{0.95}{Diagrama 8 --- Despliegue por fases (Fase 0 piloto y Fase 1/2 escalamiento).}

\subsection{8.2 Topolog\'ia de red}\label{red}

La red se organiza con cabecera en el Centro de Salud I-4 ``Altiplano'', que concentra el gabinete de telecomunicaciones (rack con servidor local, switch administrado y UPS), doble enlace WAN con \emph{failover} y firewall con segmentaci\'on por VLAN. Los puestos de salud sat\'elite I-1/I-2 operan con enlaces 4G intermitentes y, por dise\~no, en modalidad offline-first: sincronizan cuando hay enlace disponible, sin que su operaci\'on cl\'inica dependa de \'el.

\diag{d09-topologia-red.pdf}{0.92}{Diagrama 9 --- Topolog\'ia de red (cabecera I-4 y puestos sat\'elite con redundancia).}

\subsection{8.3 Centro de datos y capacidad}\label{cpd}

El dise\~no de centro de datos (Entregable N.° 9) fija el nivel de disponibilidad objetivo (clasificaci\'on Tier), el \emph{layout} f\'isico, el dimensionamiento de capacidad (c\'omputo, almacenamiento y energ\'ia) y la estrategia de virtualizaci\'on/cloud h\'ibrido para el escalamiento de Fase 1/2. En Fase 0, el ``centro de datos'' se reduce deliberadamente a un servidor \'unico contenerizado; la ruta de crecimiento hacia alta disponibilidad (r\'eplicas, balanceo, respaldo geogr\'afico) queda trazada pero no ejecutada, en coherencia con el modelo por fases. \textbf{{[}DATO PENDIENTE DE VALIDAR: nivel Tier objetivo y cifras exactas de capacidad, definidos en el Entregable N.° 9{]}}.

\section{9. Vista de seguridad}\label{seguridad}

La seguridad se estructura en capas defensivas complementarias (defensa en profundidad), desde el dispositivo hasta los datos, con una capa transversal de auditor\'ia y trazabilidad. El modelo se alinea a ISO 27005 / NIST (Entregable N.° 8) y a la Ley N.° 29733: cifrado en tr\'ansito (TLS/VPN), en reposo (AES-256) y por columna para PII; RBAC en Keycloak reforzado por \emph{Row-Level Security} en el motor; y doble auditor\'ia (FHIR \texttt{AuditEvent}/\texttt{Provenance} a nivel de aplicaci\'on m\'as \texttt{pgAudit} a nivel de motor). Un an\'alisis STRIDE resumido identifica como amenazas prioritarias la suplantaci\'on de identidad (mitigada por OIDC/MFA), la fuga de PII (mitigada por minimizaci\'on y cifrado por columna) y el repudio (mitigado por la traza inmutable).

\diag{d10-seguridad-capas.pdf}{0.9}{Diagrama 10 --- Arquitectura de seguridad por capas y flujo de identidad (OIDC/RBAC).}

\section{10. Vista de integraci\'on / interoperabilidad}\label{integracion}

La integraci\'on con el ecosistema del MINSA se concentra en el paquete \texttt{interoperabilidad\_externa}, que act\'ua como \'unico punto de contacto con los sistemas externos y a\'isla al n\'ucleo cl\'inico de sus particularidades. El backend genera \emph{bundles} FHIR idempotentes hacia el HIS-MINSA, consulta el Padr\'on Nominal para la identificaci\'on de beneficiarios y coordina con SIGA la disponibilidad de insumos; los conflictos de esquema o rechazos se marcan para revisi\'on en lugar de perderse.

\diag{d11-integracion-fhir.pdf}{0.9}{Diagrama 11 --- Arquitectura de integraci\'on FHIR con los sistemas externos.}

\section{11. Arquitectura por fases}\label{fases}

La arquitectura es deliberadamente \textbf{evolutiva}. En \textbf{Fase 0} (piloto de validaci\'on) todo el sistema corre en un servidor \'unico y el aprendizaje federado permanece inactivo aunque empaquetado. En \textbf{Fase 1/2} (escalamiento) el backend migra a un despliegue de mayor disponibilidad (r\'eplicas, balanceo, respaldo), y se activa el servidor Flower dedicado, que intercambia \'unicamente par\'ametros de modelo entre postas ---nunca datos crudos--- materializando la regla RN-09. Esta separaci\'on por fases evita sobredimensionar el piloto sin cerrar la puerta al crecimiento, y es coherente con el business case por fases del Entregable N.° 2. El Diagrama 8 (secci\'on 8.1) ilustra ambas fases.

\section{12. Integraci\'on de las tres ramas}\label{convergencia}

El valor central de este documento es evidenciar que las tres \'areas de competencia no son entregables aislados sino aportes complementarios a una \emph{misma} arquitectura. El Diagrama 12 muestra esa convergencia y la matriz siguiente la formaliza componente por componente.

\diag{d12-convergencia-ramas.pdf}{0.98}{Diagrama 12 --- Convergencia de las tres ramas (CE01, CE02, CE03) en la arquitectura integrada.}

{\small
\begin{longtable}[]{@{}
  >{\raggedright\arraybackslash}p{(\linewidth - 6\tabcolsep) * \real{0.22}}
  >{\raggedright\arraybackslash}p{(\linewidth - 6\tabcolsep) * \real{0.26}}
  >{\raggedright\arraybackslash}p{(\linewidth - 6\tabcolsep) * \real{0.26}}
  >{\raggedright\arraybackslash}p{(\linewidth - 6\tabcolsep) * \real{0.26}}@{}}
\toprule\noalign{}
Elemento de arquitectura & CE01 --- Gesti\'on de TI aporta & CE02 --- Ing. de Software aporta & CE03 --- Infraestructura aporta \\
\midrule\noalign{}
\endhead
\bottomrule\noalign{}
\endlastfoot
Drivers y alcance & Diagn\'ostico, business case, alcance por fases & Requerimientos funcionales y no funcionales & Restricciones de red, energ\'ia y capacidad \\
App m\'ovil y nodo & Proceso TO-BE y automatizaciones A1-A14 & Dise\~no de la app, inferencia y outbox & Gama de dispositivo, conectividad BLE/4G \\
Backend modular & Modelo de operaci\'on y roles & Componentes, API FHIR, l\'ogica de dominio & Dimensionamiento del servidor, contenerizaci\'on \\
Datos & Gobierno de datos y cumplimiento & Modelo relacional, PostGIS, recursos FHIR & Almacenamiento, respaldo, cifrado en reposo \\
Integraci\'on externa & Convenios y flujos institucionales & Conectores FHIR e idempotencia & Enlaces, VPN y disponibilidad de red \\
Seguridad & Pol\'iticas y cumplimiento (Ley 29733) & Authz por recurso, RLS, auditor\'ia de app & ISO 27005/NIST, cifrado, segmentaci\'on, pgAudit \\
Despliegue & Presupuesto y roadmap por fases & Artefactos desplegables y su configuraci\'on & Topolog\'ia, centro de datos, Tier y virtualizaci\'on \\
\end{longtable}
}

La lectura por filas confirma que ning\'un elemento de arquitectura queda hu\'erfano de alguna de las tres ramas, y la lectura por columnas confirma que cada \'area tiene un aporte definido y no redundante: es la evidencia de una arquitectura verdaderamente integrada.

\section{13. Riesgos arquitect\'onicos y deuda t\'ecnica}\label{riesgos}

{\small
\begin{longtable}[]{@{}
  >{\raggedright\arraybackslash}p{(\linewidth - 4\tabcolsep) * \real{0.34}}
  >{\raggedright\arraybackslash}p{(\linewidth - 4\tabcolsep) * \real{0.22}}
  >{\raggedright\arraybackslash}p{(\linewidth - 4\tabcolsep) * \real{0.44}}@{}}
\toprule\noalign{}
Riesgo arquitect\'onico & Impacto & Mitigaci\'on \\
\midrule\noalign{}
\endhead
\bottomrule\noalign{}
\endlastfoot
Servidor \'unico en Fase 0 (punto \'unico de fallo) & Alto & Respaldo diario/semanal; ruta de HA trazada para Fase 1/2; operaci\'on offline-first reduce el impacto en campo. \\
Precisi\'on cl\'inica del algoritmo sin validaci\'on de campo & Muy alto & Estrategia de tamizaje (no diagn\'ostico) + confirmaci\'on humana obligatoria; validaci\'on cl\'inica previa a escalar. \\
Par\'ametros cl\'inicos NTS 213 no disponibles a\'un & Medio & Valores parametrizados y versionados en \texttt{config} (ADR-05), no fijados en c\'odigo. \\
Deriva de conflictos de sincronizaci\'on offline & Medio & Idempotencia por UUID + estado \emph{conflicto pendiente de revisi\'on} con resoluci\'on manual auditada. \\
Barrera de entrada tecnol\'ogica baja (hardware comod\'itizado) & Medio & Ventaja sostenible en validaci\'on cl\'inica, interoperabilidad real y confianza institucional, no en el hardware. \\
Extracci\'on futura a microservicios & Bajo & Los siete paquetes ya son l\'imites de dominio; cada uno es candidato a servicio sin reescritura (ADR-02). \\
\end{longtable}
}

\begin{center}\rule{0.5\linewidth}{0.5pt}\end{center}

\begin{quote}
\textbf{Cierre.} Esta arquitectura integrada es consistente con los nueve entregables del proyecto y con sus decisiones ya tomadas; no introduce tecnolog\'ias nuevas, sino que unifica y hace expl\'icitas las relaciones entre las piezas ya dise\~nadas por las tres \'areas de competencia. Los puntos marcados \textbf{{[}DATO PENDIENTE DE VALIDAR{]}} deben cerrarse con los datos institucionales reales antes de considerar la arquitectura como definitiva para construcci\'on.
\end{quote}

\end{document}
